\section{Practical Implications (RQ5)}
\label{sec:practical-implications}

This section synthesizes the preceding results into guidance for using finite-horizon counterfactual supervision.

\textbf{When should the offline signal be trusted?} Most directly in workload/capacity regimes with high capacity-scale reuse, identifiable before closed-loop replay via a high trace-derived LRU hit rate or low all-tied fraction. In this dataset, this corresponds to metacdn- and twemcache-like regimes, where offline discriminativeness and closed-loop agreement are strongest.

\textbf{When is caution warranted?} In near-zero-signal regimes (e.g., alibaba-block/cap32), small offline regret gaps may flip signs without reflecting genuine disagreement. Caution is also needed when interpreting cold-start evaluation windows (metakv) or validation-only regimes (metacdn). A high tie rate is itself informative: it identifies where policy comparison is fundamentally unresolvable by this label family (e.g., wiki2018), preventing researchers from hunting for non-existent signal.

\textbf{Why still run closed-loop evaluation?} The offline/closed-loop association is real but bounded, resting on $n=10$ cells with at least one unexplained exception (metakv/cap128). Offline-label results are complementary diagnostics and cannot replace closed-loop performance claims.

\textbf{How should continuation-policy assumptions be interpreted?} As a tested robustness property with real limits. The labels are highly robust to MRU continuation at population scale, and acceptably robust to random continuation in the validated sample. This supports the use of canonical LRU-continuation labels in evaluated regimes, but does not imply continuation choice is universally irrelevant.
